% capitulos/02_organizacao.tex

\chapter{ORGANIZAÇÃO DIDÁTICO-PEDAGÓGICA}

\section{POLÍTICAS DE ENSINO, PESQUISA E EXTENSÃO NO ÂMBITO DO CURSO}
\subsection{Políticas de ensino}
Texto...
\subsection{Políticas de pesquisa}
Texto...
\subsection{Políticas de extensão}
Texto...

\section{ESTRUTURA CURRICULAR}

\subsection{Requisitos para integralização curricular}
Na Tabela \ref{tab:integralizacao}, é apresentada a distribuição de carga horária exigida para integralização do curso.

\begin{table}[htb]
\centering
\caption{Distribuição da carga horária exigida para integralização do curso}
\label{tab:integralizacao}
\begin{tabularx}{\textwidth}{|X|c|}
\hline
\textbf{Modalidade da Atividade} & \textbf{Carga Horária} \\ \hline
1. Componentes Curriculares Obrigatórios de Graduação & \\ \hline
\hspace{0.5cm} 1.1 Trabalho de Conclusão de Curso & \\ \hline
\hspace{0.5cm} 1.2 Estágio Curricular Obrigatório & \\ \hline
2. Componentes Curriculares Complementares de Graduação & \\ \hline
3. Atividades Complementares de Graduação & \\ \hline
4. Atividades Curriculares de Extensão & \\ \hline
\hspace{0.5cm} 4.1 Ativ. Curriculares de Extensão Vinculadas & \\ \hline
\hspace{0.5cm} 4.2 Ativ. Curriculares de Extensão Específicas & \\ \hline
\hspace{1cm} 4.2.1 Unipampa Cidadã & \\ \hline
\textbf{* Total} & \\ \hline
\end{tabularx}
\fonte{Elaborado pela Comissão do PPC (2026).}
\end{table}

\subsection{Matriz curricular}
A matriz curricular do curso é apresentada na Tabela \ref{tab:matriz}.

\begin{table}[htb]
\centering
\caption{Matriz Curricular do Curso}
\label{tab:matriz}
\resizebox{\textwidth}{!}{
\begin{tabular}{|c|c|l|c|c|c|c|c|c|c|c|}
\hline
\textbf{Sem.} & \textbf{CÓD} & \textbf{Nome do Componente} & \textbf{Pré-req.} & \textbf{Teórica} & \textbf{Prática} & \textbf{EaD Teó} & \textbf{EaD Prá} & \textbf{Ext.} & \textbf{Total} & \textbf{Créd.} \\ \hline
1º & EC001 & Algoritmos e Programação & - & 30 & 30 & 0 & 0 & 0 & 60 & 4 \\ \hline
% Adicione as demais linhas aqui
\end{tabular}
}
\fonte{Elaborado pela Comissão do PPC (2026).}
\end{table}

\subsection{Distribuição da Carga Horária do curso}
Texto...

\subsection{Abordagem dos temas transversais}
Texto...

\subsection{Flexibilização curricular}

\subsubsection{Componentes Curriculares Complementares de Graduação (CCCGs)}
\begin{table}[htb]
\centering
\caption{Componentes Curriculares Complementares de Graduação do Curso}
\label{tab:cccg}
\begin{tabular}{|c|l|c|c|c|c|c|}
\hline
\textbf{Código} & \textbf{Nome do Componente} & \textbf{Teórica} & \textbf{Prática} & \textbf{Extensão} & \textbf{Total} & \textbf{Créd.} \\ \hline
EC099 & Tópicos Especiais em IA & 30 & 30 & 0 & 60 & 4 \\ \hline
\end{tabular}
\fonte{Elaborado pela Comissão do PPC (2026).}
\end{table}

\subsubsection{Atividades Complementares de Graduação (ACG)}
Texto explicativo sobre ACGs e Tabelas dos Grupos I a IV...

\subsubsection{Mobilidade acadêmica}
Texto...

\subsubsection{Aproveitamento de estudos}
Texto...

\subsubsection{Carga horária a distância em cursos presenciais}
\paragraph{Outros recursos didáticos} Texto...
\paragraph{Equipe Multidisciplinar para Educação a Distância} Texto...

\subsection{Migração curricular e equivalências}
Texto e Tabela 8 de Migração...

\subsection{Atividades Práticas de Ensino}
Aplicável mais à saúde, mas preencher se houver laboratórios práticos intensivos...

\subsection{Estágios obrigatórios ou não obrigatórios}
Texto...

\subsection{Práticas Profissionais}
Texto...

\subsection{Trabalho de conclusão de curso}
Texto...

\subsection{Inserção da extensão no currículo do curso}
Texto...

\section{METODOLOGIAS DE ENSINO E APRENDIZAGEM}
\subsection{Interdisciplinaridade}
Texto...
\subsection{Práticas inovadoras}
Texto...
\subsection{Tecnologias de informação e comunicação (TICs)}
Texto
\subsection{Acessibilidade metodológica, curricular e pedagógica}
Texto

\section{AVALIAÇÃO DA APRENDIZAGEM}
Texto

\section{APOIO AO(À) DISCENTE}
Texto

\section{GESTÃO DO CURSO A PARTIR DO PROCESSO DE AVALIAÇÃO INTERNA E EXTERNA}
Texto